\documentclass[aspectratio=169,10pt]{beamer}
\usetheme[numbering=fraction,progressbar=frametitle]{metropolis}
\usepackage{fontspec}
\setsansfont{Latin Modern Sans}
\usepackage{amsmath,amssymb,mathtools,graphicx,booktabs}
\usepackage{polyglossia}
\setdefaultlanguage{french}
\setbeamertemplate{navigation symbols}{}
\setbeamerfont{frametitle}{size=\large}
\setbeamerfont{normal text}{size=\normalsize}
\setlength{\parskip}{.45em}
\hypersetup{pdfauthor={},pdfsubject={Lecture slides},pdfkeywords={},pdftitle={Géométrie repérée : synthèse vecteurs et droites}}
\providecommand{\N}{\mathbb{N}}
\providecommand{\Z}{\mathbb{Z}}
\providecommand{\D}{\mathbb{D}}
\providecommand{\Q}{\mathbb{Q}}
\providecommand{\R}{\mathbb{R}}
\providecommand{\C}{\mathbb{C}}
\providecommand{\vect}[1]{\overrightarrow{#1}}
\providecommand{\vi}{\vec{\imath}}
\providecommand{\vj}{\vec{\jmath}}
\providecommand{\vk}{\vec{k}}
\providecommand{\norm}[1]{\left\lVert #1\right\rVert}
\providecommand{\vc}[2]{\begin{pmatrix}#1\\#2\end{pmatrix}}
\providecommand{\deter}[2]{\det\left(#1,#2\right)}
\providecommand{\abs}[1]{\left\lvert #1\right\rvert}
\providecommand{\dd}{\,\mathrm{d}}
\providecommand{\e}{\mathrm{e}}
\providecommand{\Card}{\operatorname{Card}}
\providecommand{\E}{\operatorname{E}}
\providecommand{\V}{\operatorname{V}}
\providecommand{\degre}{^{\circ}}
\providecommand{\un}[1]{\,\mathrm{#1}}
\providecommand{\sci}[2]{#1\times10^{#2}}
\providecommand{\pgcd}{\operatorname{PGCD}}
\providecommand{\ppcm}{\operatorname{PPCM}}
\title{Géométrie repérée : synthèse vecteurs et droites}
\subtitle{Cours et exemples}
\date{}
\author{}
\begin{document}
\begin{frame}[plain]
\vfill
{\usebeamerfont{title}\inserttitle\par}
\medskip
{\usebeamerfont{subtitle}\insertsubtitle\par}
\medskip
{\usebeamercolor[fg]{progress bar}\rule{\textwidth}{.5pt}}
\vfill
\end{frame}
\begin{frame}{Parallélogramme et intersection}
\begin{columns}[T,onlytextwidth]
\begin{column}{.57\textwidth}
Dans un repère orthonormé, \(A(-4;-3)\), \(B(1;-1)\) et \(C(-1;3)\).

\begin{enumerate}
\setcounter{enumi}{0}
\item Déterminer \(D\) tel que \(ABCD\) soit un parallélogramme.

\item Déterminer son centre \(I\) de deux façons.

\item Avec \(E(2;1)\), déterminer \(K\), intersection de \((AB)\) et \((CE)\).

\item Déterminer \(F\) tel que \(\vect{AF}=\dfrac23\vect{AE}\). Les points \(D,F,B\) sont-ils alignés ?

\end{enumerate}
\end{column}
\begin{column}{.4\textwidth}
\centering
\includegraphics[width=\linewidth,height=4.5cm,keepaspectratio]{figures/seconde-12-1.pdf}
\end{column}
\end{columns}
\end{frame}
\begin{frame}{Correction : Parallélogramme et intersection}
1. \(D=A+C-B=(-6;1)\). 2. Milieu de \([AC]\) et de \([BD]\) : \(I(-5/2;0)\). 3. \((AB):2x-5y-7=0\) ; \((CE):2x+3y-7=0\). Soustraire donne \(8y=0\), puis \(x=7/2\), donc \(K(7/2;0)\). 4. \(\vect{AE}=(6;4)\), donc \(F=A+(4;8/3)=(0;-1/3)\). \(\vect{DF}=(6;-4/3)\), \(\vect{DB}=(7;-2)\) ; déterminant \(-12+28/3=-8/3\ne0\). Les trois points ne sont pas alignés.
\end{frame}
\begin{frame}{Comparer une méthode vectorielle et une méthode repérée}
Dans un triangle non aplati \(ABC\), \(I\) est le milieu de \([AB]\) et \(J\) le milieu de \([AC]\).

\begin{enumerate}
\setcounter{enumi}{0}\item Démontrer vectoriellement que \((IJ)\parallel(BC)\).
\item Vérifier avec \(A(0;0)\), \(B(6;0)\), \(C(2;4)\).
\item Pour ces coordonnées, déterminer les équations de \((IJ)\) et \((BC)\), puis le rapport \(IJ/BC\).

\end{enumerate}
\end{frame}
\begin{frame}{Correction : Comparer une méthode vectorielle et une méthode repérée}
1. \(\vect{IJ}=\vect{IA}+\vect{AJ}=\dfrac12(\vect{BA}+\vect{AC})=\dfrac12\vect{BC}\). Les vecteurs sont non nuls et colinéaires.

2. \(I(3;0)\), \(J(1;2)\) ; \(\vect{IJ}=(-2;2)\) et \(\vect{BC}=(-4;4)\).

3. \((IJ):x+y-3=0\) ; \((BC):x+y-6=0\). Leurs directeurs sont colinéaires. \(IJ=\sqrt8\), \(BC=\sqrt{32}\), donc \(IJ/BC=1/2\). La preuve vectorielle est générale ; la vérification numérique ne démontre que ce cas particulier.
\end{frame}
\begin{frame}{Comparer une méthode vectorielle et une méthode repérée}
Une figure semble montrer deux droites parallèles. Quelle différence y a-t-il entre cette observation et une preuve ?
\end{frame}
\begin{frame}{Correction : Comparer une méthode vectorielle et une méthode repérée}
La figure suggère une conjecture et peut être imprécise. Une preuve utilise des données exactes et une propriété, par exemple un déterminant nul de vecteurs directeurs non nuls. Une conclusion doit préciser les objets concernés.
\end{frame}
\begin{frame}{Reconnaître une configuration et justifier}
Dans un repère orthonormé, \(A(0;0)\), \(B(4;2)\), \(C(3;4)\) et \(D(-1;2)\).

\begin{enumerate}
\setcounter{enumi}{0}\item Montrer que \(ABCD\) est un parallélogramme.
\item Calculer les quatre longueurs et déterminer s'il s'agit d'un losange.
\item Calculer \(AC^2\), \(BD^2\) et déterminer s'il s'agit d'un rectangle.
\item Déterminer son centre et une équation de chaque diagonale.

\end{enumerate}
\end{frame}
\begin{frame}{Correction : Reconnaître une configuration et justifier}
1. \(\vect{AB}=(4;2)=\vect{DC}\) : c'est un parallélogramme. 2. \(AB=CD=\sqrt{20}=2\sqrt5\), \(BC=DA=\sqrt5\) ; les côtés adjacents ne sont pas égaux, donc pas un losange. 3. \(AC^2=3^2+4^2=25\) ; \(BD^2=(-5)^2+0^2=25\). Un parallélogramme dont les diagonales ont la même longueur est un rectangle. On peut aussi utiliser \(AB^2+BC^2=20+5=25=AC^2\) et la réciproque de Pythagore. 4. Le centre est \(I(3/2;2)\). \((AC):4x-3y=0\) ; \((BD):y=2\). Leur intersection retrouve \(I\).
\end{frame}
\end{document}
